\section{Introduction}
\label{ch:introduction}

\subsection{Background and Context}
\label{sec:background}

South Africa regulates consumer credit tightly. The \ac{NCA} \cite{sa_nca_2005} requires every
registered credit provider to assess whether an applicant can afford a proposed agreement before
granting it, and Regulation 23A of the National Credit Regulations
\cite{sa_ncr_affordability_2015} prescribes a residual-income test for doing so. Registered lenders
report their exposures to the credit bureaus, so each lender assesses an applicant against a record
of what every other registered lender has advanced.

\ac{BNPL} providers sit outside that regime. Because they levy neither interest nor fees on the
deferred amount, with late-payment penalties their only charge, they argue that they do not
transact credit within the meaning of the \ac{NCA}, and
so do not register with the \ac{NCR}, do not perform the Regulation 23A assessment, and do not
report what they have lent to any credit bureau \cite{nortonrose_bnpl_sa}. TransUnion, a consumer
credit bureau, confirms the absence of \ac{BNPL} information from South African credit reports
\cite{transunion_cps_sa_2025}. A registered lender running a lawful and diligent affordability
assessment therefore computes residual income from a liability record that is incomplete by
regulatory construction. It cannot see \ac{BNPL} obligations, and neither can any \ac{BNPL} provider
see another's, so a household can hold several concurrent facilities that no party to any of those
agreements has assessed in aggregate.

deHaan et al. \cite{dehaan2024bnpl}, tracking banking
records for 10.6 million United States consumers, find that adopting \ac{BNPL} is followed by higher
overdraft charges, credit-card interest and late fees, concentrated among liquidity-constrained
consumers. The \ac{CFPB} \cite{cfpb2025bnpl} finds that across 2021 and 2022, 63\% of \ac{BNPL}
borrowers held simultaneous loans at some point and 32\% held them across different firms. South
Africa has no equivalent evidence base, and the reporting exemption makes one impossible to
assemble from administrative data, since no party, the regulator included, can measure aggregate
\ac{BNPL} exposure.

This thesis builds an \ac{ABM} of the South African
household credit market: 5{,}000 synthetic households, constructed from national survey microdata
and subject to the statutory rules above, earn, spend, borrow, repay and default over a simulated
two-year horizon. \ac{BNPL} platforms with the observed South African contract structure are then injected
into an otherwise unchanged 2017 environment, so that any difference between the two runs is
attributable to the new lender class. The thesis contributes
three things: a survey-grounded agent-based model of South African household credit, a quantified
account of what the reporting exemption does to household distress, and a test of the
intervention repertoire regulators abroad have already adopted.

\subsection{Research Questions}
\label{sec:rqs}

\begin{quote}
    How does the exemption of \ac{BNPL} from South Africa's credit-reporting and affordability
    regime affect household credit distress, and which interventions in the \ac{BNPL} lending
    environment reduce it?
\end{quote}

Three subsidiary questions decompose it.

\begin{enumerate}
    \item Can a synthetic household population built from South African survey microdata reproduce
    credit-market behaviour it was not fitted to?
    \item Does the mutual invisibility of \ac{BNPL} platforms produce concurrent-facility
    stacking, and does population default rise with the number of platforms a household can stack
    across?
    \item Which interventions in the \ac{BNPL} lending environment reduce the resulting distress
    and default?
\end{enumerate}

Where the distress concentrates in the income distribution is treated as a finding rather than a
scope restriction: the empirical record places \ac{BNPL} harm among lower-income,
liquidity-constrained consumers \cite{hayashi2025constraints, wang2026buynow}, and results are
reported by income quintile throughout.

The first question is a precondition for the other two: a model whose households do not behave
like South African households cannot support a claim about what happens when a new lender class is
introduced to them. It is answered in two parts, by the population construction of
Section~\ref{ch:population} and by the baseline results of Section~\ref{ch:results}. The second
and third are answered in Section~\ref{ch:results}.